%%% STD - Organisation: Beginn
\begin{center}
  \makeatletter
  \textbf{
    \@title
  }
  \makeatother
  \\[1em]
  Informatik -- Eingebettete Systeme -- Hochschule Bremerhaven
\end{center}

\begin{table}[ht]
  \centering
  \begin{tabular}[t]{|p{15mm}|p{30mm}|p{30mm}|p{50mm}|r|}
    \hline
    \textbf{Version} & \textbf{Autor*inn*en} & \textbf{geändert} & \textbf{Beschreibung} & \textbf{fertig am}\\
    \hline
    0.1              & A. B.            &  Abschnitte 13 bis 16  & Erstellung des Dokumentes & 01.10.2019\\
    \hline
  \end{tabular}
  \caption{Versionshistorie}
\end{table}

\vspace*{4em}

Relevante Standards und Normen
\begin{itemize}
\item IEEE 829-2008 -- IEEE Standard for Software and System Test
  Documentation

  \cite{ieee-829-2008-system-test-documentation}
\item IEEE 1008-1987 -- IEEE Standard for Software Unit Testing

  \cite{ieee-1008-1987-unit-testing}
\item IEEE 1012-2016 -- IEEE Standard for System, Software, and
  Hardware Verification and Validation

  \cite{ieee-1012-2016-hardware-verification-validation}
\item IEEE 29119-2-2013 -- ISO/IEC/IEEE International Standard --
  Software and systems engineering -- Software testing -- Part 2:Test
  processes

  \cite{ieee-29119-2-2013-test-processes}
\item IEEE 29119-3-2013 -- ISO/IEC/IEEE International Standard --
  Software and systems engineering -- Software testing -- Part 3: Test
  documentation

  \cite{ieee-29119-3-2013-test-documentation}
\item IEEE 29119-4-2015 -- ISO/IEC/IEEE International Standard --
  Software and systems engineering -- Software testing -- Part 4: Test
  techniques

  \cite{ieee-29119-4-2013-test-techniques}
\end{itemize}
\clearpage
%%% STD - Organisation: Ende

\hinweis{Im Dokument STD werden Testfälle definiert. Diese Testfälle
  dienen der Prüfung des fertigen Produkts und prüfen die
  Anforderungen aus dem SRS-Dokument. Die gesamte Spezifikation
  beschreibt, welche Anforderungen das Produkt hat, wie man es
  herstellen und wie man es prüfen kann. Ein Testprotokoll, welches
  tatsächlich durchgeführte Prüfungen beschreibt, gehört in den
  Abschnitt Testprotokoll.}

\chapter{Einleitung}
\label{std:einleitung}

\section{Systemübersicht}
\label{std:systemuebersicht}

\section{Testverfahren}
\label{std:testverfahren}

\chapter{Testplan}
\label{std:testplan}

\section{Zu testende Produktmerkmale}
\label{std:merkmale-zu-testen}

\section{Nicht zu testende Produktmerkmale}
\label{std:merkale-nicht-zu-testen}

\section{Testwerkzeuge und Testumgebung}
\label{std:test-werkzeuge-umgebung}

\chapter{Testfälle}
\label{std:testfälle}

% Kopiere Testfallbeschreibung von Beginn bis Ende für jede Aufgabe
% Testfallbeschreibung: Beginn
\section{Testfall-N}
\label{std:testfall-n}

\subsection{Berührte Anforderungen}
\label{std:testfall-n-anforderungen-beruehrt}

\hinweis{Ein Beispiel: Dieser Testfall prüft die folgenden Anforderungen ab:}
\begin{itemize}
\item \hinweis{REQ-SYS-EXT-USER-1 Titel}
\item \hinweis{REQ-PROD-2}
\end{itemize}

\subsection{Aufgabe}
\label{std:testfall-n-aufgabe}

\subsection{Eingaben und Einstellungen}
\label{std:testfall-n-eingaben}

\subsection{Erwartete Ausgaben und Erfüllungskriterien}
\label{std:testfall-n-ausgaben-erwartet}

\subsection{Testdurchführung}
\label{std:testfall-n-durchfuehrung}

% Testfallbeschreibung: Ende

\chapter{Zusatzmaterial}
\label{std:zusatzmaterial}

\appendix
\chapter{Testprokotolle}
\label{std:appx-testprotokolle}

\section{Testergebnisse}
\label{std:testergebnisse}

\section{Fehlerberichte}
\label{std:fehlerberichte}

